\chapter{Menghitung dengan Bilangan Desimal}\label{ch:g5:decimalops}

Bilangan desimal dijumlahkan, dikurangkan dan dikalikan hampir sama
seperti bilangan bulat --- seluruh seninya adalah mengetahui di mana
titiknya jatuh. Bab ini membahas cara bersusun dan keajaiban
mengalikan atau membagi dengan $10$, $100$, $1000$. (Mengalikan dua
bilangan desimal menunggu pada \cref{ch:g6:decimals}.)

\section{Menjumlahkan dan mengurangkan}

\begin{method}[Bersusun dengan titik]\label{met:g5:decimalops:addsub}
\begin{enumerate}
  \item Tulislah bilangannya dengan \emph{\omterm{def:g4:decimals:point}{titik desimal} yang
        \omterm{def:g4:geometry:def}{sejajar}} (dengan begitu satuan berada di bawah satuan,
        persepuluhan di bawah persepuluhan);
  \item lengkapi bagian desimal yang lebih pendek dengan \omterm{def:g1:counting:zero}{nol};
  \item jumlahkan atau kurangkan seperti pada bilangan bulat;
  \item turunkan titiknya lurus ke bawah pada hasilnya.
\end{enumerate}
\end{method}

\begin{example}\label{ex:g5:decimalops:addsub}
$27.5 + 3.86$ (lengkapi $27.5 = 27.50$):
\[
\begin{array}{r}
2\,7.5\,0 \\
+\ \ \ 3.8\,6 \\
\hline
3\,1.3\,6
\end{array}
\qquad\qquad
\begin{array}{r}
1\,2.4\,0 \\
-\ \ \ 5.6\,5 \\
\hline
\ \ \,6.7\,5
\end{array}
\]
Untuk pengurangannya ($12.4 - 5.65$): lengkapi, pinjam seperti
biasa, lalu periksa dengan menjumlahkan kembali: $6.75 + 5.65 = 12.4$.
\end{example}

\begin{example}[Pelengkap di kepala]\label{ex:g5:decimalops:complement}
Berapa yang harus ditambahkan pada $7.6$ untuk mencapai $10$?
Naiklah: $7.6 \to 8$ adalah $0.4$, lalu $8 \to 10$ adalah $2$: jadi
pelengkapnya $2.4$. Berguna untuk kembalian: bayar $10$ untuk barang
seharga $7.60$, terima $2.40$ kembali.
\end{example}

\section{Sepuluh kali lebih besar, sepuluh kali lebih kecil}

\begin{proposition}[Dikali dan dibagi 10, 100, 1000]\label{prop:g5:decimalops:shift}
Mengalikan dengan $10$ membuat setiap angka bernilai sepuluh kali
lebih besar: angkanya bergeser satu tempat ke kiri --- yang terlihat
seperti \emph{titiknya berpindah satu tempat ke kanan}. Membagi
melakukan sebaliknya:
\[
3.47 \times 10 = 34.7, \qquad
3.47 \times 100 = 347, \qquad
52.1 \div 10 = 5.21, \qquad
6 \div 100 = 0.06 .
\]
\end{proposition}
\admitted

\begin{omfigure}
\begin{tikzpicture}[scale=0.9]
  \draw (0,0) grid[xstep=1.5, ystep=0.75] (6.0,2.25);
  \node[font=\small] at (0.75,2.6) {puluhan};
  \node[font=\small] at (2.25,2.6) {satuan};
  \node[font=\small] at (3.75,2.6) {persepuluhan};
  \node[font=\small] at (5.25,2.6) {perseratusan};
  \node[font=\normalsize, omExo] at (2.25,1.85) {$3$};
  \node[font=\normalsize, omExo] at (3.75,1.85) {$4$};
  \node[font=\normalsize, omExo] at (5.25,1.85) {$7$};
  \node[font=\normalsize, omDef] at (0.75,1.1) {$3$};
  \node[font=\normalsize, omDef] at (2.25,1.1) {$4$};
  \node[font=\normalsize, omDef] at (3.75,1.1) {$7$};
  \node[right, font=\small] at (6.2,1.85) {$3.47$};
  \node[right, font=\small, omDef] at (6.2,1.1)
    {$3.47 \times 10 = 34.7$};
  \draw[-Stealth, omDef, thick] (4.4,1.7) -- (3.1,1.25);
\end{tikzpicture}

{\small Mengalikan dengan $10$: setiap angka naik satu tempat ke
kiri. Titiknya sebenarnya tidak berpindah --- angkanyalah yang berpindah.}
\end{omfigure}

\begin{example}[Mengubah satuan lagi]\label{ex:g5:decimalops:convert}
Satuan ukuran berjarak pangkat sepuluh satu sama lain, jadi mengubah
satuan hanyalah pergeseran seperti ini: $2.35$ m $= 235$ cm
($\times 100$); $470$ g $= 0.47$ kg ($\div 1000$); $1.5$ L $= 150$ cL.
\end{example}

\section{Mengalikan bilangan desimal dengan bilangan bulat}

\begin{method}[Desimal dikali bilangan bulat]\label{met:g5:decimalops:mult}
Hitunglah seolah-olah tidak ada titiknya, lalu berikan pada hasilnya
sebanyak angka desimal yang dipunyai faktor desimalnya:
\[
4.35 \times 6: \quad 435 \times 6 = 2\,610
\ \to\ 4.35 \times 6 = 26.10 = 26.1 .
\]
Taksirlah agar titiknya diletakkan dengan yakin: $4.35$ dekat dengan
$4$, dan $4 \times 6 = 24$ --- jadi $26.1$, bukan $2.61$ atau $261$.
\end{method}

\begin{example}[Penjumlahan berulang tetap berlaku]\label{ex:g5:decimalops:mult}
$0.75 \times 4 = 0.75 + 0.75 + 0.75 + 0.75 = 3$: empat \omterm{def:g3:division:half}{seperempat}
menjadi satu keseluruhan, dalam pakaian desimal. Mengalikan dengan
bilangan bulat tetap berarti ``sekian kali''.
\end{example}

\begin{example}[Berbelanja]\label{ex:g5:decimalops:shopping}
Tiga buku tulis seharga $2.45$ masing-masing dan sebuah pena seharga $1.20$:
\begin{enumerate}
  \item buku tulis: $2.45 \times 3 = 7.35$;
  \item seluruhnya: $7.35 + 1.20 = 8.55$;
  \item kembalian dari $10$: $10 - 8.55 = 1.45$.
\end{enumerate}
\end{example}

\section{Latihan}

\begin{exercise}[$\star$]\label{exo:g5:decimalops:1}
Hitunglah secara bersusun: $45.7 + 8.93$; \quad $16.25 + 3.75$; \quad
$0.86 + 12.4$.
\end{exercise}

\begin{exercise}[$\star$]\label{exo:g5:decimalops:2}
Hitunglah secara bersusun lalu periksa: $9.4 - 2.72$; \quad
$20 - 13.45$; \quad $6.03 - 5.9$.
\end{exercise}

\begin{exercise}[$\star$]\label{exo:g5:decimalops:3}
Pelengkap menuju $10$ di kepala (seperti pada
\cref{ex:g5:decimalops:complement}): $6.2$; \quad $3.75$; \quad
$9.99$.
\end{exercise}

\begin{exercise}[$\star$]\label{exo:g5:decimalops:4}
Hitunglah tanpa cara bersusun:
\[
7.24 \times 10, \qquad
0.58 \times 100, \qquad
34.5 \div 10, \qquad
7 \div 100, \qquad
0.3 \times 1000 .
\]
\end{exercise}

\begin{exercise}[$\star$]\label{exo:g5:decimalops:5}
Ubahlah: $4.05$ m menjadi cm; \quad $325$ cm menjadi m; \quad $0.6$
kg menjadi g; \quad $85$ cL menjadi L.
\end{exercise}

\begin{exercise}[$\star$]\label{exo:g5:decimalops:6}
Taksirlah dulu, lalu hitunglah: $6.2 \times 4$; \quad
$3.45 \times 8$; \quad $12.5 \times 6$.
\end{exercise}

\begin{exercise}[$\star$]\label{exo:g5:decimalops:7}
Satu putaran lintasan panjangnya $0.4$ km. Berapa panjang $7$
putaran? Dan $25$ putaran?
\end{exercise}

\begin{exercise}[$\star$]\label{exo:g5:decimalops:8}
Lia membeli $2$ roti baguette seharga $1.15$ masing-masing dan
sebuah kue seharga $12.60$. Ia membayar dengan uang $20$. Hitunglah \omterm{def:g1:addition:def}{jumlah} belanjaannya dan kembaliannya.
\end{exercise}

\begin{exercise}[$\star$]\label{exo:g5:decimalops:9}
Sebuah botol berisi $1.5$ L. Berapa liter dalam satu paket berisi
$6$ botol? Sebuah keluarga minum $0.75$ L setiap hari: untuk berapa
hari satu paket itu bertahan? (Berapa kali $0.75$ muat di dalam $9$?)
\end{exercise}

\begin{exercise}[$\star\star$]\label{exo:g5:decimalops:10}
Miko menghitung $5.6 \times 3 = 15.18$, dengan alasan
``$5 \times 3 = 15$ dan $6 \times 3 = 18$''. Pakailah taksiran untuk
menunjukkan bahwa jawaban itu pasti salah, lalu hitunglah dengan benar.
\end{exercise}

\begin{exercise}[$\star\star$]\label{exo:g5:decimalops:11}
Empat putaran seorang pelari memakan waktu $65.4$ s, $64.9$ s,
$65.0$ s dan $66.1$ s.
\begin{enumerate}
  \item Hitunglah waktu seluruhnya untuk keempat putaran itu.
  \item Rekor untuk empat putaran adalah $260$ s. Apakah dia
        memecahkannya? Berapa \omterm{ex:g1:subtraction:difference}{selisihnya}?

\end{enumerate}
\end{exercise}
